\section*{ПРИЛОЖЕНИЕ Б. Расширенная методика экспериментов и угрозы валидности}

\subsection*{Б.1. Операционализация критериев успеха}

В любой системной работе по производительности есть риск подменить
цель удобными для измерения числами. Поэтому перед обсуждением
экспериментов необходимо явно сформулировать критерии успеха и
связать их с измерениями.

Для настоящей темы эти критерии таковы:

\begin{enumerate}
\item система должна уменьшить стоимость отражения одного изменения
маршрута;
\item система должна сохранить practically acceptable время чтения;
\item система должна сохранить возможность построения полного снимка;
\item система должна поддержать потоковый режим как основной
операционный сценарий.
\end{enumerate}

Только после такой фиксации становится понятно, почему \texttt{Dump},
\texttt{Lookup} и \texttt{Upsert} являются достаточным минимальным
набором метрик: первый отражает сценарий полного снимка для внешнего
клиента, второй -- сценарий точечного запроса от потребителя,
третий -- ключевую операцию streaming-режима.

\subsection*{Б.2. Репрезентативность синтетических данных}

В работе используются синтетические наборы префиксов, поскольку они
позволяют воспроизводимо варьировать размер таблицы, число записей и
тип адресного семейства. Это удобно и честно для микроизмерений, но
создаёт известную угрозу валидности: реальные BGP-таблицы имеют более
сложное распределение длин префиксов, плотности кластеризации,
уровней агрегации и перекрытий
\cite{labovitz1997instability,labovitz2000convergence}.

Поэтому результаты следует интерпретировать как:

\begin{itemize}
\item корректные для сравнительного анализа двух реализаций на общем
наборе операций;
\item достаточно сильные для подтверждения архитектурной гипотезы;
\item но не как окончательную оценку поведения на всех возможных
реальных интернет-таблицах.
\end{itemize}

В дальнейшем работу можно усилить измерениями на выгрузках RIS
\cite{ripe-ris} или RouteViews \cite{routeviews} либо на
production-дампах GoBGP в анонимизированном виде. Полезным
дополнением может оказаться использование RIS-beacons
\cite{ripe-beacons} как контрольного потока с известной семантикой.

\subsection*{Б.3. Микробенчмарки и реальная эксплуатация}

Микробенчмарки удобны, потому что дают чистое сравнение отдельных
операций. Но они не захватывают все факторы production-среды:

\begin{itemize}
\item сетевые задержки при общении с GoBGP;
\item конкурентные пользовательские запросы и подписки;
\item long-lived allocator behaviour и фрагментацию памяти;
\item влияние garbage collector на длительном uptime;
\item скачки нагрузки на фоне bursty update stream.
\end{itemize}

Это ограничение не делает эксперименты бесполезными. Оно лишь требует
аккуратной интерпретации: измерения показывают не всю систему, а ядро
её вычислительной модели.

\subsection*{Б.4. Почему сравнительный эксперимент всё равно валиден}

Несмотря на ограничения, сравнительный эксперимент остаётся сильным
по двум причинам.

Во-первых, сравниваются две реализации одного и того же abstract
contract внутри одной и той же сервисной оболочки. Все различия
концентрируются именно в backend-слое, а не в окружающей инфраструктуре.

Во-вторых, именно та операция, которая является целевой для новой
архитектуры~-- \texttt{Upsert}~-- демонстрирует радикальный выигрыш.
Положительный результат возникает не в случайной второстепенной
метрике, а в главной.

\subsection*{Б.5. Шум и повторяемость измерений}

Даже на одном и том же железе microbenchmarks подвержены
флуктуациям. На них влияют фоновая активность системы, планировщик
потоков, состояние кэшей CPU, поведение памяти, тепловой режим и
множество других факторов. Поэтому важно:

\begin{itemize}
\item использовать повторные прогоны (\texttt{count=2} в работе);
\item интерпретировать близкие значения осторожно;
\item не делать сильных выводов из различий, лежащих в пределах
шума.
\end{itemize}

Именно поэтому выводы по shard sweep в работе формулируются как
качественные, а не как утверждение о том, что «16 шардов всегда
лучше 8» или наоборот.

\subsection*{Б.6. Валидность результатов по IPv6}

Особого внимания заслуживают результаты IPv6. Здесь разница между
backend-вариантами оказывается особенно заметной. Однако и здесь
необходимо осторожное чтение.

С одной стороны, ухудшение IPv6 lookup в текущей реализации~--
реальный факт, который нельзя скрывать. С другой стороны, оно
отражает конкретный выбор алгоритма LPM поверх структуры хранения, а
не приговор самому подходу «stream + ART». Иначе говоря, измеряется
не абстрактный ART вообще, а конкретная реализация longest prefix
match над ним.

\subsection*{Б.7. Конструктивные угрозы валидности}

Основные угрозы валидности можно сгруппировать так.

\textbf{Construct validity:}

\begin{itemize}
\item измеряемые операции упрощают реальную эксплуатационную картину;
\item часть сценариев использует синтетические наборы префиксов;
\item нагрузка пользователей и нагрузка update stream пока не
моделируются одновременно.
\end{itemize}

\textbf{Internal validity:}

\begin{itemize}
\item различия между вариантами могли бы быть искажены деталями
реализации, а не только структурой данных;
\item некоторые инженерные тюнинги могли быть выполнены несимметрично
(например, ART сравнивается без сортировки дампа, но с ней также есть
прогон).
\end{itemize}

\textbf{External validity:}

\begin{itemize}
\item результаты не следует напрямую переносить на специализированные
forwarding-plane устройства;
\item результаты не гарантируют те же относительные величины на
другой runtime-платформе или другом объёме таблиц.
\end{itemize}

\subsection*{Б.8. Почему эти угрозы не разрушают главный вывод}

Главный вывод работы устойчив к перечисленным угрозам. Даже если
абсолютные цифры изменятся на другом железе или на другом наборе
префиксов, качественная разница между:

\begin{itemize}
\item полной перестройкой состояния на каждое изменение и
\item локальным изменением одного шарда
\end{itemize}

сохраняется. А именно эта разница и является ядром исследуемого
архитектурного перехода.

\subsection*{Б.9. Что стоит добавить в следующем цикле исследований}

Для усиления работы и последующих публикаций можно предложить
следующие направления развития экспериментальной базы:

\begin{enumerate}
\item измерения на реальных снимках RIS \cite{ripe-ris} и RouteViews
\cite{routeviews};
\item долговременные replay-прогоны update streams (часы, дни);
\item смешанная нагрузка \texttt{Lookup}~+~\texttt{Subscribe}~+~updates;
\item измерения времени восстановления после потери соединения и
ресинхронизации;
\item профилирование GC и memory usage на длинных интервалах
(\texttt{pprof});
\item сравнение с дополнительными структурами LPM, например
Poptrie-like baseline \cite{asai2015poptrie} или вариант,
построенный на Tree Bitmap \cite{eatherton2004treebitmap};
\item методическое сравнение с подходами в духе controlled prefix
expansion \cite{srinivasan1999prefixexpansion}.
\end{enumerate}

%==============================================================
